\documentclass[11pt]{article}

\usepackage[margin=1in]{geometry}
\usepackage[T1]{fontenc}
\usepackage{amsmath,amssymb}
\usepackage{microtype}
\usepackage{xcolor}
\usepackage{enumitem}
\usepackage{booktabs}
\usepackage{array}
\usepackage{titlesec}
\usepackage{listings}
\usepackage[hidelinks]{hyperref}

\definecolor{accent}{RGB}{32,89,120}
\definecolor{muted}{RGB}{93,99,106}
\definecolor{codebg}{RGB}{247,248,250}
\definecolor{codeframe}{RGB}{220,224,230}

\titleformat{\section}
  {\large\bfseries\color{accent}}
  {\thesection}{0.6em}{}
\titleformat{\subsection}
  {\normalsize\bfseries}
  {\thesubsection}{0.6em}{}

\setlist[itemize]{leftmargin=1.4em, itemsep=0.18em, topsep=0.25em}
\setlist[enumerate]{leftmargin=1.6em, itemsep=0.18em, topsep=0.25em}
\setlength{\parskip}{0.55em}
\setlength{\parindent}{0pt}

\lstdefinestyle{python}{
  language=Python,
  basicstyle=\ttfamily\small,
  backgroundcolor=\color{codebg},
  frame=single,
  rulecolor=\color{codeframe},
  showstringspaces=false,
  columns=fullflexible,
  keepspaces=true,
  breaklines=true,
  tabsize=4,
  keywordstyle=\color{accent}\bfseries,
  commentstyle=\color{muted},
  stringstyle=\color{teal!70!black}
}

\newcommand{\R}{\mathbb{R}}
\newcommand{\vx}{\mathbf{x}}
\newcommand{\vu}{\mathbf{u}}
\newcommand{\vz}{\mathbf{z}}
\newcommand{\vphi}{\boldsymbol{\phi}}
\newcommand{\vm}{\mathbf{m}}
\newcommand{\vn}{\mathbf{n}}
\newcommand{\vI}{\mathbf{I}}
\newcommand{\vw}{\mathbf{w}}
\newcommand{\vr}{\mathbf{r}}
\newcommand{\dt}{\Delta t}

\title{\vspace{-1.2cm}Low-rank Connectivity and the Dynamics It Produces}
\author{A short note on the framework used by Dubreuil et al. (2022)}
\date{}

\begin{document}
\maketitle
\vspace{-0.8cm}

\section{The starting point: a recurrent neural network}

The paper studies firing-rate recurrent neural networks. A common continuous-time
form is
\begin{equation}
  \tau \frac{d x_i}{dt}
  =
  -x_i
  +
  \sum_{j=1}^{N} J_{ij}\,\phi(x_j)
  +
  \sum_{s=1}^{N_{\mathrm{in}}} I_i^{(s)} u_s(t),
  \label{eq:rnn_component}
\end{equation}
where:
\begin{itemize}
  \item $x_i(t)$ is the activation of neuron $i$;
  \item $\phi(x_i)$ is its firing rate after a static nonlinearity;
  \item $J_{ij}$ is the recurrent weight from neuron $j$ to neuron $i$;
  \item $u_s(t)$ are external inputs, such as stimulus or context signals;
  \item $I_i^{(s)}$ is the input weight from input channel $s$ to neuron $i$.
\end{itemize}

In vector form,
\begin{equation}
  \tau \dot{\vx}
  =
  -\vx + J\phi(\vx) + \sum_s \vI^{(s)}u_s(t).
\end{equation}
A completely general $N \times N$ recurrent matrix $J$ can have rank as large as
$N$. Such a matrix can create many independent recurrent feedback directions.
The low-rank assumption sharply restricts this.

\section{What ``low rank'' means}

A rank-$R$ recurrent connectivity matrix can be written as a sum of $R$ outer
products:
\begin{equation}
  J
  =
  \frac{1}{N}
  \sum_{r=1}^{R} \vm^{(r)} \left(\vn^{(r)}\right)^{\top},
  \label{eq:lowrank_matrix}
\end{equation}
or componentwise,
\begin{equation}
  J_{ij}
  =
  \frac{1}{N}
  \sum_{r=1}^{R} m_i^{(r)} n_j^{(r)}.
\end{equation}

Each term has two population-length vectors:
\begin{itemize}
  \item $\vn^{(r)}$ is a \emph{read-in} or \emph{selection} vector. It defines
  which population activity pattern is measured.
  \item $\vm^{(r)}$ is a \emph{write-out} vector. It defines which pattern of
  recurrent input is sent back into the neurons.
\end{itemize}

For rank one,
\begin{equation}
  J = \frac{1}{N}\vm\vn^\top,
\end{equation}
so recurrent processing has a simple form:
\begin{equation}
  J\phi(\vx)
  =
  \vm
  \left(
    \frac{1}{N}\vn^\top \phi(\vx)
  \right).
\end{equation}
The network first projects the current firing rates onto $\vn$, producing one
scalar, and then feeds that scalar back along $\vm$.

\section{The latent variable}

The scalar produced by the projection is the key low-dimensional variable:
\begin{equation}
  \kappa(t)
  =
  \frac{1}{N}
  \sum_{j=1}^{N} n_j \phi(x_j(t)).
  \label{eq:kappa_rank_one}
\end{equation}
For rank one, the recurrent input to neuron $i$ is simply
\begin{equation}
  m_i \kappa(t).
\end{equation}
Thus the full $N$-dimensional network is not using an arbitrary recurrent input
pattern. Its recurrent input always lies along the single population direction
$\vm$. For rank $R$, there are $R$ such latent variables:
\begin{equation}
  \kappa_r(t)
  =
  \frac{1}{N}
  \sum_{j=1}^{N} n_j^{(r)} \phi(x_j(t)),
  \qquad r=1,\ldots,R.
\end{equation}
The recurrent input is then
\begin{equation}
  J\phi(\vx)
  =
  \sum_{r=1}^{R} \vm^{(r)}\kappa_r(t).
  \label{eq:recurrent_low_dim}
\end{equation}

This is the main mathematical reason low-rank networks are interpretable: the
large neural population can be described through a small set of recurrent latent
variables $\kappa_1,\ldots,\kappa_R$.

\section{Why low-rank connectivity creates low-dimensional dynamics}

Insert the low-rank form into the RNN:
\begin{equation}
  \tau \dot{\vx}
  =
  -\vx
  +
  \sum_{r=1}^{R} \vm^{(r)}\kappa_r(t)
  +
  \sum_s \vI^{(s)}u_s(t).
  \label{eq:low_dim_state}
\end{equation}
This equation says that changes in the neural state are built only from:
\begin{itemize}
  \item the recurrent directions $\vm^{(1)},\ldots,\vm^{(R)}$;
  \item the input directions $\vI^{(1)},\ldots,\vI^{(N_{\mathrm{in}})}$;
  \item the leak term $-\vx$.
\end{itemize}

If the initial state is in the subspace spanned by the recurrent and input
directions, it stays there. Even when the full network has thousands of neurons,
the meaningful trajectories can live in a much smaller subspace. This is what the
paper means by reducing the network to latent dynamics.

\section{A rank-one example}

With one input $u(t)$ and rank-one recurrent connectivity,
\begin{equation}
  \tau \dot{x}_i
  =
  -x_i + m_i \kappa(t) + I_i u(t),
\end{equation}
\begin{equation}
  \kappa(t)
  =
  \frac{1}{N}
  \sum_j n_j \phi(x_j(t)).
\end{equation}

If $\phi$ were linear, the latent dynamics would be approximately
\begin{equation}
  \tau \dot{\kappa}
  =
  -\kappa
  +
  \sigma_{nm}\kappa
  +
  \sigma_{nI}u(t),
  \label{eq:linear_kappa}
\end{equation}
where
\begin{equation}
  \sigma_{nm} = \frac{1}{N}\sum_i n_i m_i,
  \qquad
  \sigma_{nI} = \frac{1}{N}\sum_i n_i I_i.
\end{equation}

These overlaps control the computation:
\begin{itemize}
  \item $\sigma_{nI}$ says how strongly the input drives the latent variable;
  \item $\sigma_{nm}$ says how strongly the latent variable feeds back onto itself.
\end{itemize}
If feedback is weak, $\kappa$ tracks the input briefly. If feedback is tuned close
to the leak, $\kappa$ can integrate information over time. If the nonlinearity is
strong and feedback is positive, the same low-dimensional system can support
attractor-like decision states.

\section{Where population structure enters}

The paper then asks: what statistics of the neuron-specific parameters are needed
for a task? For each neuron, collect its coordinates in connectivity space:
\begin{equation}
  \left(
    m_i^{(1)},\ldots,m_i^{(R)},
    n_i^{(1)},\ldots,n_i^{(R)},
    I_i^{(1)},\ldots,I_i^{(N_{\mathrm{in}})},
    w_i
  \right).
\end{equation}
Across neurons, these coordinates form a distribution.

\subsection*{Single random population}

In a fully random population structure, these neuron coordinates are modeled as
draws from one joint distribution, often a multivariate Gaussian:
\begin{equation}
  \mathbf{a}_i \sim \mathcal{N}(0,\Sigma).
\end{equation}
The effective latent dynamics are governed by global covariances such as
$\sigma_{nm}$ and $\sigma_{nI}$. There are no distinct groups of neurons with
different computational roles.

\subsection*{Multiple subpopulations}

In a non-random population structure, the coordinates are better described as a
mixture:
\begin{equation}
  p(\mathbf{a})
  =
  \sum_{p=1}^{P}
  \alpha_p \mathcal{N}(\mu_p,\Sigma_p).
\end{equation}
Each component corresponds to a subpopulation. The subpopulations can have
different covariances, so the same input may modulate one group strongly and
another weakly.

In the nonlinear case, this is crucial. Dubreuil et al. write the effective couplings
as population-weighted quantities of the form
\begin{equation}
  \widetilde{\sigma}_{ab}
  =
  \sum_p
  \alpha_p
  \sigma_{ab}^{(p)}
  \langle \phi' \rangle_p.
\end{equation}
Here $\langle \phi' \rangle_p$ is the gain of subpopulation $p$. Inputs can change
these gains differently across subpopulations. This lets context reshape the
effective latent dynamics without necessarily driving the same latent variable
directly.

\section{The computational message}

Low rank controls the \emph{dimension} of the recurrent dynamics. A rank-$R$
network has $R$ main recurrent latent variables.

Population structure controls how inputs, feedback, and readout are arranged
across neurons. A single random population can implement many tasks by using
low-dimensional dynamics. But tasks requiring flexible input-output remapping can
require multiple subpopulations, because different groups need to have different
gain modulations and different effective roles.

In one sentence:

\begin{quote}
Low-rank connectivity says the network computes through a small number of
collective variables; non-random population structure says different neurons are
organized so that inputs can flexibly reshape the dynamics of those collective
variables.
\end{quote}

\section*{Glossary}

\begin{tabular}{@{}ll@{}}
\toprule
Term & Meaning \\
\midrule
$J$ & recurrent connectivity matrix \\
rank $R$ & number of independent recurrent feedback directions \\
$\vn^{(r)}$ & vector that reads population activity into latent variable $r$ \\
$\vm^{(r)}$ & vector that writes latent variable $r$ back into the population \\
$\kappa_r$ & low-dimensional recurrent latent variable \\
$\vI^{(s)}$ & input vector for input channel $s$ \\
$\vw$ & readout vector for producing the network output \\
population structure & distribution of neuron coordinates in connectivity space \\
\bottomrule
\end{tabular}

\clearpage
\section{Additional primer for implementing the assignment}

The original note above explains why low-rank connectivity is useful. The
remaining sections add the implementation details needed to build and train the
specific rank-one RNN from the assignment.

\section{The concrete model in the task}

The task asks for a recurrent neural network with one-dimensional input
$u(t) \in \R$, one-dimensional output $z(t) \in \R$, and $N$ recurrent units.
The continuous-time state of neuron $i$ is $x_i(t)$. Its firing rate is
\begin{equation}
  r_i(t) = \phi(x_i(t)) = \tanh(x_i(t)).
\end{equation}
The dynamics are
\begin{equation}
  \tau \frac{d x_i}{dt}
  =
  -x_i
  +
  \sum_{j=1}^{N} J_{ij}\phi(x_j)
  +
  I_i u(t),
  \qquad i=1,\ldots,N,
  \label{eq:task_continuous}
\end{equation}
with neuronal time constant $\tau = 100\text{ ms}$.

The recurrent matrix is rank one:
\begin{equation}
  J = \frac{1}{N}\vm\vn^\top,
  \qquad
  J_{ij} = \frac{1}{N}m_i n_j.
  \label{eq:rank_one_J}
\end{equation}
The output is a linear readout of the rates:
\begin{equation}
  z(t)
  =
  \frac{1}{N}\sum_{i=1}^{N}w_i\phi(x_i(t)).
  \label{eq:task_output}
\end{equation}

Only $\vm$ and $\vn$ are trainable. The input vector $\vI$ and output vector
$\vw$ are fixed after initialization. All trainable and fixed parameters are
sampled from standard normal distributions, except $\vw$, which is sampled with
standard deviation $4$.

\subsection{Vector form}

Let $\vx(t) \in \R^N$, $\vI,\vm,\vn,\vw \in \R^N$, and
$\phi(\vx)$ mean applying $\tanh$ elementwise. Then
\begin{equation}
  \tau \dot{\vx}
  =
  -\vx
  +
  \frac{1}{N}\vm\vn^\top\phi(\vx)
  +
  \vI u(t).
  \label{eq:vector_continuous}
\end{equation}
Because $J$ is rank one, it is better not to explicitly build the full
$N \times N$ matrix. Define the scalar recurrent coordinate
\begin{equation}
  \kappa(t)
  =
  \frac{1}{N}\vn^\top \phi(\vx(t)).
  \label{eq:kappa}
\end{equation}
Then the recurrent input is simply
\begin{equation}
  J\phi(\vx(t)) = \vm \kappa(t).
\end{equation}
So the implemented update can use one dot product with $\vn$, then multiply by
$\vm$. This is both clearer and cheaper than forming $J$.

\section{What an RNN is doing}

A recurrent neural network is a dynamical system with memory. A feedforward
network computes
\begin{equation}
  y = f_\theta(u)
\end{equation}
from the current input alone. An RNN computes from both the current input and an
internal state:
\begin{equation}
  \vx_{t+1} = F_\theta(\vx_t, u_t),
  \qquad
  z_t = G_\theta(\vx_t).
\end{equation}
The state $\vx_t$ carries information from earlier inputs. In this task, the
state is a population of firing-rate units. The leak term $-\vx$ makes old
activity decay, the recurrent term feeds current activity back into the network,
and the input term injects the stimulus.

The nonlinearity matters. Here $\phi = \tanh$ maps activations to rates in
$(-1,1)$. Around zero, $\tanh(x) \approx x$, so the network behaves almost
linearly. For large positive or negative activations, $\tanh$ saturates and
$\phi'(x) = 1 - \tanh^2(x)$ becomes small. Saturation can stabilize activity, but
it can also make gradients small during training.

\section{From the differential equation to code}

The task specifies forward Euler simulation with
\[
  \dt = 20\text{ ms}.
\]
Since $\tau = 100\text{ ms}$,
\[
  \alpha = \frac{\dt}{\tau} = \frac{20}{100} = 0.2.
\]
Forward Euler replaces $\dot{\vx}(t)$ by
\[
  \frac{\vx_{t+1} - \vx_t}{\dt}.
\]
Substituting into Equation~\ref{eq:vector_continuous} gives
\begin{equation}
  \vx_{t+1}
  =
  \vx_t
  +
  \alpha
  \left[
    -\vx_t
    +
    \vm\kappa_t
    +
    \vI u_t
  \right],
  \label{eq:euler_update}
\end{equation}
where
\begin{equation}
  \kappa_t = \frac{1}{N}\vn^\top\tanh(\vx_t).
\end{equation}
Equivalently,
\begin{equation}
  \vx_{t+1}
  =
  (1-\alpha)\vx_t
  +
  \alpha\vm\kappa_t
  +
  \alpha\vI u_t.
  \label{eq:euler_expanded}
\end{equation}
This is the exact recurrence to implement.

If the stimulus lasts $75$ time steps, the real duration is
\[
  75 \times 20\text{ ms} = 1500\text{ ms}.
\]

\section{Tensor shapes for a batch implementation}

Use a batch dimension so many trials can be simulated at once. One convenient
layout is:
\begin{center}
\begin{tabular}{@{}lll@{}}
\toprule
Quantity & Shape & Meaning \\
\midrule
\texttt{u} & \texttt{(batch, time)} & scalar input sequence for each trial \\
\texttt{x} & \texttt{(batch, N)} & current recurrent state \\
\texttt{rates} & \texttt{(batch, N)} & \texttt{tanh(x)} \\
\texttt{m, n, I, w} & \texttt{(N,)} & population vectors \\
\texttt{kappa} & \texttt{(batch,)} & recurrent latent coordinate \\
\texttt{z} & \texttt{(batch,)} & scalar output at one time step \\
\texttt{outputs} & \texttt{(batch, time)} & full output sequence \\
\bottomrule
\end{tabular}
\end{center}

With this layout, the low-rank computations are:
\begin{align}
  \kappa_t^{(b)}
  &=
  \frac{1}{N}\sum_{i=1}^{N} n_i \tanh(x_{t,i}^{(b)}),
  \\
  z_t^{(b)}
  &=
  \frac{1}{N}\sum_{i=1}^{N} w_i \tanh(x_{t,i}^{(b)}),
\end{align}
where $b$ indexes the trial in the batch.

In PyTorch, these become:
\begin{lstlisting}[style=python]
rates = torch.tanh(x)              # (batch, N)
kappa = (rates * n).mean(dim=1)    # (batch,)
z = (rates * w).mean(dim=1)        # (batch,)
\end{lstlisting}
The use of \texttt{mean(dim=1)} implements the factor $1/N$.

\section{Initialization and the role of the output scale}

The vectors should be initialized as
\begin{align}
  m_i &\sim \mathcal{N}(0,1), &
  n_i &\sim \mathcal{N}(0,1), &
  I_i &\sim \mathcal{N}(0,1), &
  w_i &\sim \mathcal{N}(0,4^2).
\end{align}
The larger standard deviation for $\vw$ compensates for the averaging in the
readout. At random initialization, if the rates are only weakly aligned with
$\vw$, the sum
\[
  \frac{1}{N}\sum_i w_i r_i
\]
has a typical scale proportional to $\operatorname{std}(w)/\sqrt{N}$ rather than
$\operatorname{std}(w)$. A larger output scale helps keep the initial output from
being extremely small for moderate $N$. It does not make $\vw$ trainable; it only
sets the fixed readout scale.

\section{Trainable parameters versus fixed buffers}

In PyTorch, trainable parameters should be \texttt{nn.Parameter} objects. Fixed
tensors that belong to the model should be registered as buffers. Buffers move
with the model when calling \texttt{model.to(device)}, are saved in the
\texttt{state\_dict}, but are not updated by the optimizer.

For this task:
\begin{center}
\begin{tabular}{@{}lll@{}}
\toprule
Vector & PyTorch representation & Optimizer updates it? \\
\midrule
$\vm$ & \texttt{nn.Parameter} & yes \\
$\vn$ & \texttt{nn.Parameter} & yes \\
$\vI$ & \texttt{register\_buffer("I", ...)} & no \\
$\vw$ & \texttt{register\_buffer("w", ...)} & no \\
\bottomrule
\end{tabular}
\end{center}

\section{A minimal PyTorch module}

This is the core implementation pattern. It avoids forming $J$ and unrolls the
dynamics over time.

\begin{lstlisting}[style=python]
import torch
from torch import nn


class RankOneRNN(nn.Module):
    def __init__(self, n_units: int, dt_ms: float = 20.0, tau_ms: float = 100.0):
        super().__init__()
        self.n_units = n_units
        self.alpha = dt_ms / tau_ms

        self.m = nn.Parameter(torch.randn(n_units))
        self.n = nn.Parameter(torch.randn(n_units))

        self.register_buffer("I", torch.randn(n_units))
        self.register_buffer("w", 4.0 * torch.randn(n_units))

    def forward(self, u: torch.Tensor, x0: torch.Tensor | None = None):
        # u has shape (batch, time)
        batch_size, n_steps = u.shape

        if x0 is None:
            x = u.new_zeros(batch_size, self.n_units)
        else:
            x = x0

        outputs = []
        states = []

        for t in range(n_steps):
            rates = torch.tanh(x)

            z = (rates * self.w).mean(dim=1)
            outputs.append(z)
            states.append(x)

            kappa = (rates * self.n).mean(dim=1)
            recurrent = kappa[:, None] * self.m[None, :]
            external = u[:, t, None] * self.I[None, :]

            x = x + self.alpha * (-x + recurrent + external)

        outputs = torch.stack(outputs, dim=1)
        states = torch.stack(states, dim=1)
        return outputs, states
\end{lstlisting}

Important shape details:
\begin{itemize}
  \item \texttt{kappa} has shape \texttt{(batch,)}.
  \item \texttt{kappa[:, None]} has shape \texttt{(batch, 1)}.
  \item \texttt{self.m[None, :]} has shape \texttt{(1, N)}.
  \item Their product broadcasts to \texttt{(batch, N)}.
\end{itemize}

\section{Training objective}

Training means choosing $\vm$ and $\vn$ so the output sequence $z_t$ matches a
target sequence $y_t$. For the assignment, use a network of size
\[
  N = 128.
\]
The target is only defined during the final decision epoch. If the simulated
sequence has length $T_{\mathrm{sim}}$, define the decision window as the last
\[
  T_{\mathrm{dec}} = 15
\]
time steps. The loss should not penalize the network during earlier stimulus or
delay time steps. It only asks whether the final decision-period output is close
to the target.

For a mini-batch of $B$ trials, the assignment's loss is
\begin{equation}
  \mathcal{L}(\vm,\vn)
  =
  \frac{1}{B T_{\mathrm{dec}}}
  \sum_{b=1}^{B}
  \sum_{s=1}^{T_{\mathrm{dec}}}
  \left(
    z_{t_s}^{(b)} - y_{t_s}^{(b)}
  \right)^2,
  \label{eq:decision_loss}
\end{equation}
where $t_s = T_{\mathrm{sim}} - T_{\mathrm{dec}} + s$ indexes the final
15 time steps. The problem statement writes the trial index as $i$; here $b$ is
used instead to avoid confusing trials with neurons.

The instruction says to use stochastic gradient descent, specifically Adam, with
\[
  B = 32,
  \qquad
  \eta = 5 \times 10^{-3}.
\]
In practice, this means each parameter update uses one mini-batch of 32 trials,
computes the average decision-window MSE, backpropagates that scalar loss, and
lets Adam update only $\vm$ and $\vn$. A successful implementation should reach a
loss below roughly $5 \times 10^{-2}$, possibly after about 1000 parameter
updates.

In code:
\begin{lstlisting}[style=python]
decision_steps = 15

outputs, states = model(u_batch)           # outputs: (batch, time)
output_decision = outputs[:, -decision_steps:]

# If target is already a time series, keep only its decision window.
# If target is one scalar per trial, repeat it across the decision window.
if target_batch.ndim == 2:
    target_decision = target_batch[:, -decision_steps:]
else:
    target_decision = target_batch[:, None].expand_as(output_decision)

loss = ((output_decision - target_decision) ** 2).mean()
\end{lstlisting}

\section{Task-specific training loop}

A full training step has four conceptual stages:
\begin{enumerate}
  \item Generate or sample a mini-batch of input trials and targets.
  \item Run the RNN forward through the full input sequence, storing outputs and
  hidden states.
  \item Compute MSE only on the last 15 decision time steps.
  \item Backpropagate through the unrolled recurrence and update $\vm,\vn$ with
  Adam.
\end{enumerate}

\begin{lstlisting}[style=python]
model = RankOneRNN(n_units=128)
optimizer = torch.optim.Adam(model.parameters(), lr=5e-3)

batch_size = 32
decision_steps = 15

for update in range(1000):
    u_batch, target_batch = make_batch(batch_size)

    optimizer.zero_grad(set_to_none=True)
    outputs, states = model(u_batch)

    output_decision = outputs[:, -decision_steps:]
    if target_batch.ndim == 2:
        target_decision = target_batch[:, -decision_steps:]
    else:
        target_decision = target_batch[:, None].expand_as(output_decision)

    loss = ((output_decision - target_decision) ** 2).mean()
    loss.backward()
    optimizer.step()

    if update % 100 == 0:
        print(update, float(loss))
\end{lstlisting}

This code assumes \texttt{make\_batch(32)} returns an input tensor with shape
\texttt{(32, time)}. The target can either be a full target time series with
shape \texttt{(32, time)} or one desired decision value per trial with shape
\texttt{(32,)}.

\section{How training works: backpropagation through time}

The forward pass unrolls the recurrence:
\begin{align}
  \vx_1 &= F_\theta(\vx_0, u_0),\\
  \vx_2 &= F_\theta(\vx_1, u_1),\\
  &\vdots\\
  \vx_T &= F_\theta(\vx_{T-1}, u_{T-1}),
\end{align}
where $\theta = \{\vm,\vn\}$ in this task.

The loss depends on all states through the outputs:
\[
  \mathcal{L}
  =
  \sum_{t \in \text{decision window}} \ell(G(\vx_t), y_t).
\]
This is important: the error is measured only at the end, but the hidden state at
the end depends on every earlier input and every earlier recurrent update.
Backpropagation through time is ordinary backpropagation applied to the unrolled
computation graph. The same parameters $\vm$ and $\vn$ are reused at every time
step, so their gradients accumulate contributions from every step:
\[
  \frac{\partial \mathcal{L}}{\partial \theta}
  =
  \sum_{t \in \text{decision window}}
  \frac{\partial \mathcal{L}}{\partial \vx_t}
  \frac{\partial \vx_t}{\partial \theta}.
\]
The term $\partial \vx_t / \partial \theta$ already contains the whole history
from the initial state up to time $t$. The dependence of later states on earlier
states introduces products of Jacobians:
\[
  \frac{\partial \vx_t}{\partial \vx_s}
  =
  \prod_{q=s}^{t-1}
  \frac{\partial F_\theta(\vx_q,u_q)}{\partial \vx_q},
  \qquad s < t.
\]
If these products shrink, gradients vanish. If they grow, gradients explode.
The leak, the factor $\alpha=0.2$, the low rank, and the saturation of $\tanh$
all affect this gradient flow.

PyTorch builds the graph automatically while the loop runs. Calling
\texttt{loss.backward()} computes the gradients; calling
\texttt{optimizer.step()} updates only \texttt{nn.Parameter} tensors.
Do not detach \texttt{x} or \texttt{states} inside the forward loop during normal
training; detaching would cut the gradient path from the decision loss back to
earlier stimulus processing. Gradient clipping is not part of the mathematical
model, but it can be useful if training produces very large gradients.

\section{Debugging failed convergence}

If the loss does not get below roughly $5 \times 10^{-2}$, first debug the
training setup before interpreting the result as a scientific failure.
\begin{enumerate}
  \item \textbf{Overfit one fixed mini-batch.} Generate one batch of 32 trials,
  reuse exactly that batch for every update, and check whether the loss can be
  driven very low. If this fails, suspect a bug in the recurrence, target
  alignment, loss window, or optimizer setup.
  \item \textbf{Train a full-rank network.} Temporarily replace the rank-one
  recurrent term by a trainable full matrix. If the full-rank model learns but
  the rank-one model does not, the data and loss pipeline are probably working.
  If the full-rank model also fails, the issue is likely in the implementation
  rather than in the rank constraint.
  \item \textbf{Inspect dynamics after convergence.} Plot representative inputs,
  outputs, targets, and hidden-state summaries. Useful hidden summaries include
  $\kappa_t$, the readout $z_t$, and a few hidden units or low-dimensional
  projections of \texttt{states}.
\end{enumerate}

The most common target-alignment bug is accidentally computing loss over the
whole sequence. That trains the network to match undefined or irrelevant targets
during stimulus presentation. For this assignment, slice the last 15 time steps:
\begin{lstlisting}[style=python]
output_decision = outputs[:, -15:]
target_decision = target[:, -15:]  # or label[:, None].expand_as(output_decision)
loss = ((output_decision - target_decision) ** 2).mean()
\end{lstlisting}

\section{What gradients hit in this rank-one model}

The trainable vectors influence the network differently.

The vector $\vn$ determines what population activity is measured:
\[
  \kappa_t = \frac{1}{N}\sum_i n_i r_{t,i}.
\]
Changing $\vn$ changes the scalar feedback signal extracted from the rates.

The vector $\vm$ determines where that scalar is written back:
\[
  \text{recurrent input}_{t,i} = m_i\kappa_t.
\]
Changing $\vm$ changes which neurons receive positive or negative recurrent
drive.

Because $\vw$ is fixed, training cannot directly rotate the output readout.
Instead, it must shape the recurrent dynamics so that the fixed readout
\[
  z_t = \frac{1}{N}\vw^\top \tanh(\vx_t)
\]
becomes task-correct.

\section{Low-rank interpretation}

A full recurrent matrix has $N^2$ entries and can mix neural activity in many
directions. The rank-one matrix
\[
  J = \frac{1}{N}\vm\vn^\top
\]
has only two length-$N$ vectors. It performs two operations:
\begin{enumerate}
  \item read the current rates along $\vn$ to get one scalar $\kappa$;
  \item write that scalar back into the population along $\vm$.
\end{enumerate}
This makes the recurrent part of the dynamics low-dimensional:
\[
  \tau\dot{\vx} = -\vx + \vm\kappa(t) + \vI u(t).
\]
Even though the state $\vx$ has $N$ components, the recurrent feedback has only
one degree of freedom. The main computation can often be understood by tracking
$\kappa(t)$ and the output $z(t)$.

For intuition, suppose $\tanh(x) \approx x$. Then
\[
  \tau \dot{\kappa}
  \approx
  -\kappa
  +
  \sigma_{nm}\kappa
  +
  \sigma_{nI}u(t),
\]
where
\[
  \sigma_{nm} = \frac{1}{N}\sum_i n_i m_i,
  \qquad
  \sigma_{nI} = \frac{1}{N}\sum_i n_i I_i.
\]
The overlap $\sigma_{nI}$ controls how strongly the input drives the latent
coordinate. The overlap $\sigma_{nm}$ controls recurrent feedback. Training
$\vm$ and $\vn$ changes these overlaps and thereby changes whether activity
decays quickly, integrates evidence, or moves toward stable decision-like states.

\section{Practical implementation checks}

Before training on the full task, test the model in small pieces:
\begin{enumerate}
  \item Check that \texttt{model(u)} returns \texttt{outputs} with shape
  \texttt{(batch, time)} and \texttt{states} with shape
  \texttt{(batch, time, N)}.
  \item Check that only \texttt{m} and \texttt{n} appear in
  \texttt{list(model.parameters())}.
  \item Check that \texttt{I} and \texttt{w} appear in
  \texttt{model.state\_dict()} but not in the optimizer parameters.
  \item Feed zero input and verify that the zero initial state stays zero.
  \item Feed a simple constant input and verify that states and outputs are
  finite, not \texttt{nan} or \texttt{inf}.
  \item Run one training step and verify that \texttt{m.grad} and
  \texttt{n.grad} are not \texttt{None}.
\end{enumerate}

Example checks:
\begin{lstlisting}[style=python]
u = torch.randn(8, 75)
model = RankOneRNN(n_units=128)
outputs, states = model(u)

assert outputs.shape == (8, 75)
assert states.shape == (8, 75, 128)
assert torch.isfinite(outputs).all()

loss = outputs.pow(2).mean()
loss.backward()
assert model.m.grad is not None
assert model.n.grad is not None
assert model.I.grad is None
assert model.w.grad is None
\end{lstlisting}

\section{Common pitfalls}

\begin{itemize}
  \item Forgetting the factor $1/N$. In code, use \texttt{mean(dim=1)} or divide
  sums by \texttt{n\_units}.
  \item Accidentally making $\vI$ or $\vw$ trainable. They should be buffers for
  this task.
  \item Building the full $J$ matrix unnecessarily. For rank one, use
  \texttt{kappa * m}.
  \item Mixing time and batch dimensions. Pick one layout and keep it consistent.
  \item Detaching the state inside the time loop during normal full-sequence
  training. Detaching cuts the gradient path through time.
  \item Using in-place tensor operations on values needed for autograd. Prefer
  assignments like \texttt{x = x + ...}.
  \item Computing the output only after the update when the target expects output
  before the update, or vice versa. Choose the convention deliberately and keep
  targets aligned with it.
\end{itemize}

\section{Minimal mental model}

At each time step, the network does the following:
\begin{enumerate}
  \item Convert activations to firing rates:
  \[
    \vr_t = \tanh(\vx_t).
  \]
  \item Read the population along $\vn$:
  \[
    \kappa_t = \frac{1}{N}\vn^\top\vr_t.
  \]
  \item Produce the output through the fixed readout:
  \[
    z_t = \frac{1}{N}\vw^\top\vr_t.
  \]
  \item Combine leak, recurrent feedback, and input:
  \[
    \vx_{t+1}
    =
    \vx_t
    +
    0.2
    \left[
      -\vx_t + \vm\kappa_t + \vI u_t
    \right].
  \]
  \item During training, compare $z_t$ to the target, backpropagate through the
  unrolled steps, and update only $\vm$ and $\vn$.
\end{enumerate}

That is the core of the assignment: implement this recurrence correctly,
generate outputs over time, define a task loss, and let PyTorch's autograd tune
the two low-rank connectivity vectors.

\section*{Glossary}

\begin{tabular}{@{}ll@{}}
\toprule
Term & Meaning \\
\midrule
$\vx_t$ & recurrent state or activation vector at time step $t$ \\
$\tanh(\vx_t)$ & firing-rate vector after the nonlinearity \\
$\tau$ & neuronal time constant; here $100\text{ ms}$ \\
$\dt$ & Euler step size; here $20\text{ ms}$ \\
$\alpha$ & discretization factor $\dt/\tau = 0.2$ \\
$J$ & recurrent connectivity matrix \\
rank one & one recurrent read-in/write-out direction pair \\
$\vn$ & trainable vector that reads rates into $\kappa$ \\
$\vm$ & trainable vector that writes $\kappa$ back into neurons \\
$\kappa_t$ & scalar recurrent latent coordinate \\
$\vI$ & fixed input vector \\
$\vw$ & fixed output readout vector \\
$z_t$ & scalar network output \\
BPTT & backpropagation through time; backprop on the unrolled recurrence \\
\bottomrule
\end{tabular}

\end{document}
